\documentclass{mathos}

\begin{document}
\renewcommand\bibname{Literatura}


\begin{titlepage}
\begin{center}
    
\includegraphics[keepaspectratio=true,scale=0.4]{MathosLogo.png} \\ \vspace{1mm}
       \textsc{Sveučilište Josipa Jurja Strossmayera u Osijeku} \\
    \vspace{2mm}
    \textsc{Odjel za matematiku}
\end{center}

\begin{center}
    \vspace{4mm}
    {Sveučilišni preddiplomski studij Matematika i računarstvo}
\end{center}

\vspace{25mm}
\begin{center}
    {\LARGE{\bf UI sustav za dinamičko povećanje rezolucije i generiranje umjetnih okvira u interaktivnoj grafici}}

    \vspace{15mm}
    {\large{\sc Završni rad}}
\end{center}

\vspace{50mm}

\begin{minipage}[t]{0.47\textwidth}
	{Mentor:} \normalsize\vspace{3mm}
	% Titula prevedena s "Associate Professor" na mathos.unios.hr profilu;
	% nije nađen hrvatski oblik na stranici, provjeri prije predaje.
	{\bf\\ \large{izv.\,prof.\,dr.\,sc. Domagoj Ševerdija}}
\end{minipage}
\hfill
\begin{minipage}[t]{0.47\textwidth}\raggedleft
	{{Kandidati:}{\normalsize\vspace{3mm} \bf\\ \large{Ante Galić}\\ \large{Ivo Janković}}}
\end{minipage}

\vspace*{\fill}
\centering{\normalsize{Osijek, 2026.}}
\end{titlepage}


\tableofcontents
\thispagestyle{empty}


\chapter{Uvod}\label{Uvod}
\pagenumbering{arabic}

Razlučivost i frekvencija osvježavanja zaslona rastu brže od rasterizacijske
snage grafičkih procesora. Posljedica je da igre i interaktivne 3D aplikacije
sve češće ne renderiraju svaki piksel svakog prikazanog okvira izravno, nego
rekonstruiraju veći dio slike iz manjeg broja stvarno izračunatih uzoraka —
kroz prostor (\emph{upscaling}) i kroz vrijeme (\emph{generiranje
interpoliranih okvira}). AMD-ov FidelityFX Super Resolution 3 (FSR3) jedan je
od nekoliko komercijalnih sustava koji te dvije tehnike kombinira u jedan
cjevovod; ovaj rad ne koristi FSR3 kao gotovu biblioteku, nego iz opisa
njegovih tehnika (dostupnih kao javna dokumentacija, ne kao izvorni kod)
iznova implementira ključne dijelove: temporalni upscaler, procjenu gibanja
iz slike (\emph{optical flow}) i generiranje međuokvira, u vlastitom
cjevovodu pisanom od nule u C++20 i OpenGL 4.6 compute shaderima.

Ključna postavka koja cijeli rad čini mjerljivim, a ne dojmljivim, je vlastiti
renderer testne scene: kako se scena i kamera u cijelosti kontroliraju, moguće
je renderirati i \emph{ground truth} — puni native okvir kao referencu za
upscaler, i pravi među-okvir na trenutku $t-0,5\Delta t$ kao referencu za
interpolaciju. Kvaliteta rezultata je time PSNR/SSIM brojka naspram poznate
istine, ne subjektivan dojam iz igranja.

Rad je podijeljen u dvanaest poglavlja. Poglavlje~\ref{TeorijskaPodloga}
postavlja teorijsku podlogu (uzorkovanje i aliasing, temporalna akumulacija,
procjena gibanja, reprojekcija i disokluzija). Poglavlje~\ref{PregledRjesenja}
daje pregled postojećih komercijalnih i istraživačkih rješenja i smješta ovaj
rad među njih. Poglavlje~\ref{Arhitektura} opisuje arhitekturu vlastitog
sustava. Poglavlja~\ref{ImplUpscaler}--\ref{ImplFrameGen} opisuju
implementaciju upscalera, procjene gibanja i generiranja okvira, pass po
pass. Poglavlje~\ref{NaucenaKomponenta} opisuje naučenu komponentu — malu
mrežu koja uči birati između kandidata koje klasični dio cjevovoda već
proizvodi. Poglavlje~\ref{Metodologija} opisuje metodologiju mjerenja, s
naglaskom na ground-truth referencu. Poglavlje~\ref{Rezultati} donosi
rezultate: kvalitetu, performanse, latenciju, ablacije i provjeru na dva
različita GPU-a. Poglavlje~\ref{Rasprava} raspravlja ograničenja i
odstupanja od izvornog plana rada, a Poglavlje~\ref{Zakljucak} zaključuje i
predlaže moguće nastavke.


\chapter{Teorijska podloga}\label{TeorijskaPodloga}

\section{Uzorkovanje i aliasing}

Rasterizacija procjenjuje integral svjetlosti koji piksel prima jednim ili
malim brojem uzoraka po pikselu. Prema teoremu uzorkovanja, signal koji
sadrži frekvencije iznad polovice frekvencije uzorkovanja (Nyquistove
frekvencije) ne može se iz uzoraka jednoznačno rekonstruirati — visoke
frekvencije se \emph{preklope} (alias) u niske i pojave se kao vidljivi
artefakti: nazubljeni rubovi, treperenje tankih detalja, moiré uzorci na
teksturama visoke frekvencije. Rješenje nije bolji filtar nakon uzorkovanja,
nego \textbf{više uzoraka prije} njega — supersampling, ili, ekonomičnije,
uzorci raspoređeni tako da se njihov prosjek kroz vrijeme približi gušćem
uzorkovanju.

Ovo je razlog zašto renderer u ovom radu, kad treba \emph{referencu} za
mjerenje kvalitete, ne koristi jednostruki render u prikaznoj razlučivosti,
nego $2\times2$ supersampling usrednjen u linearnom prostoru boje prije
tone mappinga (Poglavlje~\ref{Metodologija}): bodovanje naspram aliasirane
referencije bi sustavno nagrađivalo zamućenje, jer zamućeni filtar
``pogodi'' aliasiranu referencu bliže u srednjekvadratnoj pogrešci nego
filtar koji ispravno rekonstruira rub.

\section{Temporalna akumulacija}

Ako se kamera svaki okvir pomakne za pod-pikselni razmak (\emph{jitter}) i
rezultati se akumuliraju kroz vrijeme s odgovarajućim ponderiranjem, niz
okvira zajedno djeluje kao gušće uzorkovanje jednog okvira — temporalni
antialiasing (TAA). Karis \cite{Karis} formalizira ovaj pristup za igre u
stvarnom vremenu: svaki okvir doprinosi jedan novi uzorak po pikselu na
novoj sub-pikselnoj fazi, a povijest (prethodni akumulirani rezultat) donosi
sve prijašnje. Niz faza uzorkovanja obično se bira iz Haltonove sekvence
(baza 2, 3), jer nizak-discrepancy niz jednoliko popunjava jedinični kvadrat
brže od slučajnog uzorkovanja.

Cijena je da akumulacija zahtijeva \textbf{reprojekciju}: povijest je
snimljena za prošlu poziciju kamere (i prošle položaje pomičnih objekata), pa
se prije miješanja s trenutnim okvirom mora preslikati na trenutnu geometriju
pomoću vektora gibanja. Yang i sur. \cite{YangSurvey} u svom pregledu
tehnika temporalnog antialiasinga sustavno klasificiraju izvore pogreške
ovog koraka: pogrešni ili nepostojeći vektori gibanja (sjene, refleksije,
promjene sjenčanja nisu opisane geometrijskim gibanjem), \textbf{disokluzija}
(površina koja u prošlom okviru nije bila vidljiva nema valjanu povijest) i
\textbf{ghosting} (povijest koja se pogrešno prihvati kao valjana kad
zapravo ne pripada trenutnoj površini). Sve tri kategorije izravno odgovaraju
mehanizmima opisanima u Poglavlju~\ref{ImplUpscaler} (depth clip za
disokluziju, color-box clamp za ghosting, dilatacija vektora za rubove
objekata).

Temporalna akumulacija je, po konstrukciji, niskopropusni filtar u vremenu:
prigušuje visoke frekvencije koje se ne poklapaju savršeno između okvira.
To je poanta kad je riječ o aliasingu, ali postaje mana kad se scena stvarno
mijenja brže nego što akumulacija stigne pratiti — analitički izvod tog
efekta (formula~\eqref{eq_nyquist}) jedan je od glavnih nalaza
Poglavlja~\ref{ImplUpscaler}.

\section{Procjena gibanja i optical flow}

Vektor gibanja dobiven iz rasterizacije (razlika projicirane pozicije istog
vrha geometrije između dva okvira) postoji samo za ono što se doista crta
trokutima, s poznatom transformacijom. Optical flow rješava opći problem:
procijeniti polje gibanja \textbf{izravno iz dvije slike}, bez pristupa
geometriji koja ih je proizvela.

Klasična formulacija (Fleet i Weiss \cite{FleetWeiss}) polazi od
pretpostavke \textbf{konstantnosti svjetline}: sjaj točke na površini se ne
mijenja između okvira, samo se pomiče,
\begin{equation}
I(x, y, t) = I(x + u\,\delta t,\; y + v\,\delta t,\; t + \delta t),
\label{eq_brightness_constancy}
\end{equation}
gdje je $(u,v)$ traženi vektor gibanja. Linearizacijom
\eqref{eq_brightness_constancy} Taylorovim redom dobiva se jednadžba
optičkog toka $I_x u + I_y v + I_t = 0$ — jedna linearna jednadžba s dvije
nepoznanice po pikselu, dakle problem je \textbf{nedovoljno određen}
(\emph{ill-posed}). Dvije klasične posljedice te nedovoljne određenosti:

\begin{itemize}
\item \textbf{Aperture problem.} Unutar homogene ravne plohe ili duž ruba
koji je lokalno pravac, gradijent $I_x, I_y$ ne određuje komponentu gibanja
duž tog ruba — sve vrijednosti te komponente daju isti lokalni odgovor.
Rješava se dodatnim pretpostavkama (glatkoća polja, veći prozor za
promatranje) ili prihvaćanjem da procjena u takvim područjima nasljeđuje
grublju/susjednu procjenu.
\item \textbf{Loša postavljenost problema} u područjima niskog kontrasta ili
periodične strukture — više različitih vektora gibanja daje jednako dobro
podudaranje (isti minimum kriterija), pa je izbor osjetljiv na šum.
\end{itemize}

Ove dvije pojave nisu apstraktne u ovom radu: izravno objašnjavaju zašto
implementirani procjenitelj (Poglavlje~\ref{ImplOpticalFlow}) treba kaznu
udaljenosti od predikcije (rješava lošu postavljenost) i zašto sustavno
griješi na periodičnim teksturama i ravnim plohama (aperture problem), a
mjereno je koliko griješi, ne samo zašto.

Praktične implementacije za grafiku u stvarnom vremenu, uključujući
\verb+FfxOpticalFlow+ koji je uzor implementaciji u ovom radu, ne rješavaju
linearizirani sustav diferencijalno (Lucas--Kanade), nego \textbf{diskretnom
pretragom} (block matching): za blok piksela luminancije traži se pomak koji
minimizira kriterij razlike (npr. sumu apsolutnih razlika) unutar prozora
pretrage, u piramidi razlučivosti \emph{coarse-to-fine} da bi se pokrili i
veliki pomaci. Ovo je aproksimacija koja žrtvuje sub-pikselnu preciznost i
teorijsku elegantnost gradijentnih metoda za robusnost na velike, brze
pomake tipične za igre — kompromis koji je u ovom radu izmjeren, ne samo
naveden (Poglavlje~\ref{ImplOpticalFlow}).

\section{Reprojekcija i disokluzija}

Reprojekcija — dohvaćanje vrijednosti iz prošlog okvira na poziciji koju
vektor gibanja pokazuje — pretpostavlja da površina koja je tamo bila i dalje
postoji i vidljiva je. Dvije pretpostavke mogu propasti: povijest s te
pozicije nikad nije zapisana za \emph{ovu} površinu (nova površina se otkriva
— disokluzija), ili je zapisana za drugu, sada prekrivenu površinu (obrnuta
disokluzija). Standardan test je \textbf{provjera dubine}: usporedba dubine
koju je površina imala prošli okvir (poznata ili rekonstruirana) s dubinom
koju prošli dubinski spremnik ima na reprojiciranoj poziciji — neslaganje
iznad praga znači da povijest ne pripada ovoj površini i mora se odbaciti ili
umanjiti njezin utjecaj. Ovo je princip iza \emph{depth clip} mehanizma
opisanog u Poglavlju~\ref{ImplUpscaler} i iza dviju maski disokluzije u
generiranju međuokvira (Poglavlje~\ref{ImplFrameGen}).


\chapter{Pregled postojećih rješenja}\label{PregledRjesenja}

Ovo poglavlje smješta rad među postojeća komercijalna i istraživačka
rješenja za upscaling i generiranje okvira. Detaljnije bilješke s izvorima po
tvrdnji (URL-ovi, razlikovanje potvrđenih i nepotvrđenih navoda) čuvaju se
odvojeno od ovog rada; ovdje je sažetak i, na kraju poglavlja, izravna veza
sa svakim implementiranim modulom.

\section{NVIDIA DLSS}

DLSS~2 (2020.) \cite{DLSSNvidia} uvodi naučenu temporalnu rekonstrukciju: konvolucijski
autoenkoder na Tensor jedinicama prima boju niske razlučivosti, igrine
vektore gibanja i prethodni izlaz visoke razlučivosti, treniran naspram
$16K$ offline-renderirane referentne slike. DLSS~3 (2022.) dodaje generiranje
okvira kao zaseban dio: ulazi su dva okvira, hardverski optical flow (Ada
Optical Flow Accelerator) i igrini podaci; zahtijeva RTX~40. DLSS~4 (2025.)
mijenja mrežu za super-rezoluciju s CNN-a na transformer i seli generiranje
okvira s hardverskog optical flowa na naučeni model. DLSS~4.5 (2026.) dodaje
drugu generaciju transformera i FP8 na RTX~40/50.

Reflex~2 Frame Warp \cite{ReflexNvidia} je iznimka među NVIDIA-inim rješenjima: kasna
reprojekcija (ekstrapolacija, ne interpolacija) koja iz zadnjeg unosa i
zadnjeg renderiranog okvira izračuna novu kameru i warpa okvir neposredno
prije prikaza.

\section{AMD FSR}

FSR~1 je čisto prostorni, klasičan pristup (EASU + RCAS) — ista dva
mehanizma implementirana u ovom radu (Poglavlje~\ref{ImplUpscaler}). FSR~2
je temporalan i klasičan. \textbf{FSR~3/3.1} \cite{FSR3GPUOpen} \textbf{je klasičan i u dijelu za
generiranje okvira}: optical flow (temelj AMD Fluid Motion Frames) radi po
blokovima $8\times8$ bez preklapanja, piramida s SAD pretragom — struktura
koju M6 ovog rada slijedi, s gušćim poljem (vektor po pikselu, ne po bloku).
FSR~4 (2025.) je naučeni upscaler na FP8 matričnim jedinicama RDNA~4; FSR
``Redstone'' (2025.) dodaje naučeno generiranje okvira koje se miješa s
klasičnom reprojekcijom vektora. Za Polaris (RX~400/500, hardver za koji je
ovaj rad izvorno planiran) AMD nema izjavu — vrijede samo klasični FSR~1/2/3.

\section{Intel XeSS}

XeSS~2 (2024.) \cite{XeSSIntel} je komercijalno rješenje najbliže M9 ovog rada: tri stupnja —
reprojekcija igrinim vektorima, naučeni model za reprojekciju optičkim
tokom, i \textbf{drugi naučeni model s funkcijom mješavine koji bira između
oba kandidata}. To je gotovo doslovno arhitektura M9-a (mreža ne crta
piksele, bira/miješa kandidate), samo s naučenim, ne klasičnim, optičkim
tokom kao jednim od kandidata, i na hardveru s XMX jedinicama.

\section{Apple, Sony, Arm, Qualcomm}

Apple MetalFX ima temporalni naučeni upscaler (od 2022.) i interpolator
okvira (od 2025.). Sony PSSR (PlayStation~5 Pro, 2024.) je naučeni upscaler
na namjenskom ML hardveru; nadograđeni PSSR (2026.) dolazi iz zajedničkog
razvoja s AMD-om i koristi isti model kao FSR Upscaling, treniran na drugim
podacima, u INT8. \textbf{Arm Neural Super Sampling} \cite{ArmNSS} (najavljen 2025.) je
konceptualno najbliži analog M9 pristupa u cijelom pregledu: UNet predviđa
filterske jezgre i temporalne koeficijente, a klasični akumulator ih
primjenjuje — mreža opet ne crta piksele nego parametrizira klasični dio.
Qualcommov Snapdragon GSR je klasičan (bez ML-a).

\section{Istraživački modeli interpolacije okvira (VFI)}

RIFE, FILM, IFRNet i EMA-VFI (2022.--2023.) rade isključivo iz RGB parova
okvira, bez igrinih vektora, dubine ili odvajanja korisničkog sučelja —
moraju procijeniti cijelo gibanje samo iz piksela. Cijena je desetci do
stotine milisekundi na vrhunskim GPU-ima (RIFE: 16~ms na 480p na TITAN~X
Pascal), naspram budžeta od $\approx$8~ms za cijeli skok 60$\to$120~fps.
RIFE ima 9,8~milijuna parametara — 1361$\times$ više od mreže u ovom radu
(7199 parametara). Trenirani su na prirodnom videu, ne na aliasiranim,
jitteriranim renderima s korisničkim sučeljem, i nemaju integraciju s
pacingom ili Reflex-stilskim smanjenjem latencije — zato se ne koriste
izravno u igrama.

\section{Presjek i veza s ovim radom}

Zajednički obrazac komercijalnih rješenja: inferencija je gusto matrično
množenje u 8-bitnim formatima (FP8 na RTX~40/50 i RDNA~4, INT8 na RDNA~3
portu i na PS5~Pro, XMX/DP4a na Intelu). Bez namjenske 8-bitne podrške isti
model košta znatno više — točno razlog zašto FSR~4 na RDNA~3 mora biti
prekvantiziran u INT8, i točno razlog zašto poglavlje~\ref{NaucenaKomponenta}
ovog rada mjeri granicu male naučene mreže \textbf{bez} takvog hardvera. Svi
komercijalni sustavi za generiranje okvira \textbf{interpoliraju} (drže
zadnji stvarni okvir, pa idu uz Reflex/Anti-Lag~2/XeLL za latenciju) — isto
radi i ovaj rad (Poglavlje~\ref{ImplFrameGen}), uz izmjerenih $+8$--$16$~ms
latencije. Svi odvajaju korisničko sučelje od scene — isto radi M8 ovog
rada.

\begin{table}[hbtp]
\centering
\small
\begin{tabular}{|l|l|l|l|}\hline
Rješenje & Proizvođač & Vrsta & Klasično / naučeno \\\hline
DLSS 3/4 FG & NVIDIA & FG & naučena sinteza, (bivše) HW optical flow \\\hline
FSR 3/3.1 & AMD & SR+FG & klasično; block-SAD optical flow $8\times8$ \\\hline
FSR Redstone FG & AMD & FG & naučeno gibanje+izgled + reprojekcija \\\hline
XeSS 2 FG & Intel & FG & reprojekcija + naučeni tok + naučena mješavina \\\hline
Arm NSS & Arm & SR & naučen UNet, jezgre za klasični akumulator \\\hline
RIFE/FILM/... & istraživanje & VFI & naučen tok, samo iz RGB parova \\\hline
\textbf{ovaj rad} & --- & SR+FG & klasično (M3--M8) + naučena mješavina (M9) \\\hline
\end{tabular}
\caption{Sažetak usporedbe; puna tablica sa svim izvorima u pratećoj
dokumentaciji rada.}
\label{table_ch3_usporedba}
\end{table}


\chapter{Arhitektura vlastitog sustava}\label{Arhitektura}

\section{Odluke i njihovo obrazloženje}

\begin{table}[hbtp]
\centering
\begin{tabular}{|l|l|}\hline
Odluka & Izbor \\\hline
Jezik / API & C++20 + OpenGL 4.6 compute shaderi (GLSL 460) \\\hline
Scena & vlastiti mali renderer + glTF (Sponza, Damaged Helmet) \\\hline
Referentne implementacije & čitaju se kao opis algoritma, ne kopiraju \\\hline
Naučeni dio & mali modul + poglavlje analize, ne zamjena jezgre \\\hline
\end{tabular}
\caption{Ključne odluke plana rada.}
\label{table_odluke}
\end{table}

OpenGL 4.6 compute shaderi izabrani su nad Vulkanom jer sadrže sve potrebne
značajke (\verb+imageAtomicMax+ nad \verb+r32ui+ formatom, nužno za scatter
polja vektora u generiranju okvira; GPU timestamp queries; storage buffere)
uz znatno manje infrastrukturnog koda — vrijeme rada ide u algoritam, ne u
rukovanje sinkronizacijom i memorijskim barijerama koje Vulkan traži
eksplicitno. Vlastiti mali renderer (ne gotov engine) je nužan, ne stilski
izbor: bez izravne kontrole nad jitterom kamere, dubinskim spremnikom i
vektorima gibanja temporalni upscaling se ne može ni implementirati ni
mjeriti — gotov engine bi te podatke sakrio iza vlastite abstrakcije.
Referentna FSR3 dokumentacija (javno dostupni API-neutralni opisi tehnika)
čita se kao pseudokod; implementacija u ovom radu je vlastita, iz opisa, a
ne prijepis izvornog koda FidelityFX SDK-a.

\section{Cjevovod}

\begin{center}
\begin{verbatim}
       ulaz (misc/tipkovnica)
              |
              v
  Modul A: test-bed renderer
  kamera (jitter Halton 2,3) -> G-buffer @ render rezolucija
  boja HDR | dubina (reverse-Z) | vektori gibanja | reactive maska
              |
       +------+------+
       |             |
       v             v
  Modul B:       Modul C:
  temporalni     optical flow
  upscaler       (piramidalni
  (6 passeva)    block matching)
       |             |
       | upscaled(t-1), upscaled(t)   OF vektori
       +------+------+
              v
  Modul D: generiranje okvira (9 passeva)
  -> interpolirani okvir na t - 0,5
              |
              v
  UI kompozicija -> frame pacing (nit prikaza) -> present
              |
              v
  Modul F: mjerenje (metrike, timing)
\end{verbatim}
\end{center}

Modul B (Poglavlje~\ref{ImplUpscaler}) i Modul C (Poglavlje~\ref{ImplOpticalFlow})
rade paralelno na istim ulazima iz Modula A i hrane Modul D
(Poglavlje~\ref{ImplFrameGen}), koji rezultat generiranja predaje redu prikaza.
Modul E, naučena komponenta (Poglavlje~\ref{NaucenaKomponenta}), djeluje unutar
Modula D, kao zamjena za ručno napisano pravilo miješanja kandidata u
prolazima 10--12. Modul F (Poglavlje~\ref{Metodologija}) nije zaseban dio
cjevovoda nego instrumentacija koja prati sve ostale module.

\section{Ground truth: ključna prednost postavke}

Renderer u ovom radu, jer je vlastiti, može renderirati dva oblika
\emph{ground trutha} koje gotov engine ne izlaže: \textbf{puni native okvir}
u prikaznoj razlučivosti bez jittera (referenca za upscaler) i \textbf{pravi
okvir na trenutku $t-0,5\Delta t$}, dobiven pomicanjem scene i kamere na
taj trenutak i renderiranjem bez akumulirane povijesti (referenca za
generiranje okvira). Oba se snimaju supersamplirano ($2\times2$) i usrednjena
u linearnom prostoru boje prije tone mappinga (obrazloženje u
Poglavlju~\ref{TeorijskaPodloga}, odjeljak o uzorkovanju).

Ova postavka je razlog zašto se kvaliteta u ovom radu izvještava kao PSNR/SSIM
brojka naspram poznate istine, umjesto kao subjektivan dojam — i razlog zašto
su ablacije u poglavljima~\ref{ImplUpscaler}--\ref{NaucenaKomponenta}
mjerljive, a ne samo opisane.

\section{Razvojni hardver}

Rad je razvijan na GPU-u AMD Radeon RX~580 (Polaris, Mesa upravljački
program), izabranom jer nema namjenske matrične/tenzorske jedinice — okolnost
koja izravno motivira mjerenje granice naučene komponente u
Poglavlju~\ref{NaucenaKomponenta}. Kompletno mjerenje ponovljeno je i na
dva novija GPU-a (NVIDIA RTX~5070 i AMD RX~7800~XT), s rezultatima u
odjeljku~\ref{ProvjeraHardver}, kao provjera da zaključci ne ovise o jednom
uređaju.


\chapter{Implementacija upscalera}\label{ImplUpscaler}

Upscaler se u ovom radu razvijao u tri koraka, svaki mjerljiv zasebno: prostorni
baseline (bez ikakve vremenske komponente), minimalni temporalni upscaler i,
naposljetku, puni upscaler u stilu AMD FSR3. Svaki sljedeći korak mora
nadmašiti prethodni na istoj sceni i istoj metrici, inače nema opravdanja za
dodanu složenost.

\section{Prostorni baseline (nearest, bilinear, bicubic, FSR1)}

Prva stepenica ne pokušava biti dobra — svrha joj je postaviti donju granicu
ispod koje temporalni pristupi ne smiju pasti, i pokazati zašto je vremenska
komponenta uopće potrebna. Implementirana su četiri prostorna filtra:
najbliži susjed i bilinearno uzorkovanje (hardverski filtri GPU-a), bicubic
(Catmull-Rom, odnosno Mitchell-Netravali s parametrima $B=0$, $C=1/2$) i AMD
FSR1 (EASU + RCAS).

\textbf{FSR1} je prostorni upscaler bez povijesti. EASU (\emph{Edge-Adaptive
Spatial Upsampling}) prolaz koristi 12 tapova (4$\times$4 bez kutova): četiri
preklapajuća petotapna križa, težinski po bilinearnim težinama kvadranta, daju
procjenu smjera i jačine ruba; kernel se zatim rotira u smjeru ruba i rasteže
okomito na njega, s prozorskom funkcijom koja ga ograničava i deringing
korakom koji izlaz stegne u raspon četiri najbliža tapa. RCAS (\emph{Robust
Contrast Adaptive Sharpening}) koji slijedi računa 3$\times$3 križ, po kanalu
min/max prsten, i ograničava negativnu laticu preostalim prostorom do crne i
do bijele boje, skaliranu faktorom $2^{-s}$ gdje je $s$ parametar oštrine.

Obje implementacije su vlastite, izvedene iz objavljenog opisa algoritma
(\verb+ffx_fsr1.h+ se čita kao referenca, ne kopira — vidi princip u odjeljku
o referentnim implementacijama).

Bitna implementacijska odluka je \textbf{redoslijed tone mappinga i
upscalera}: FSR1 je definiran nad display-referred vrijednostima u
$[0,1]$ — RCAS-ov limiter pretpostavlja da je vrh skale 1, a EASU-ova detekcija
rubova pretpostavlja perceptivni razmak između vrijednosti. Renderer radi u
linearnom HDR-u, pa lanac ide:
\begin{center}
linearni HDR (G-buffer) $\to$ tone mapping $\to$ upscaler (EASU $\to$ RCAS) $\to$ prikaz
\end{center}
Suprotan redoslijed bi EASU-u dao pogrešne smjerove rubova i RCAS koji ne
oštri ništa.

\paragraph{Rezultati.} Mjereno na sceni Sponza, 1920$\times$1080, 120 okvira,
skriptirana kamera, referenca supersamplirana $2\times2$ i usrednjena u
linearnom HDR-u prije tone mappinga (razlog je objašnjen u
Poglavlju~\ref{Metodologija}):

\begin{table}[hbtp]
\centering
\begin{tabular}{|l|c|c|}\hline
Upscaler & Quality 1,5$\times$ (PSNR / SSIM) & Performance 2,0$\times$ (PSNR / SSIM) \\\hline
Nearest & 32,15 dB / 0,8966 & 31,59 dB / 0,8794 \\\hline
Bilinear & 35,25 dB / 0,9333 & 33,11 dB / 0,8895 \\\hline
Bicubic & 35,37 dB / 0,9420 & 33,13 dB / 0,8997 \\\hline
FSR1 (EASU+RCAS) & 33,09 dB / 0,9117 & 31,92 dB / 0,8785 \\\hline
EASU sam (RCAS isključen) & \textbf{35,42 dB} / 0,9406 & --- \\\hline
\end{tabular}
\caption{Prostorni upscaleri, Sponza 1080p.}
\label{table_m3_rezultati}
\end{table}

Zanimljiv je nalaz da je \textbf{RCAS oštrina monotono štetna} na ovoj sceni:
ablacija parametra oštrine od najjačeg (0) do praktički isključenog (8) daje
i PSNR i SSIM koji rastu monotono kako se oštrenje gasi (30,89~dB pri 0 do
35,42~dB pri 8). Optimum je, dakle, degenerirani slučaj u kojem RCAS ne radi
ništa — EASU je već izvukao svu informaciju koju jedan okvir sadrži, pa je
oštrina ovdje perceptivni parametar, ne parametar vjernosti, i tako ju treba
i izvještavati (ablacijom, ne jednom brojkom). Ovo je izravna posljedica
scene: Sponza ima mipmapirane teksture s malo aliasinga za popraviti.

Po cijeni, FSR1 na ovoj sceni \textbf{ne ubrzava ništa}: cijena upscalera je
fiksna i vezana za izlaznu razlučivost, dok je ušteda vezana za ulaznu —
Sponza se na testnom GPU-u renderira native u $\approx$1~ms na 1080p, pa
ušteda od nižeg ulaznog uzorkovanja ne pokriva cijenu EASU+RCAS-a i tone
mappinga. Ovo je stvarno ograničenje mjerne postavke (scena je jednostavnija
od realnog opterećenja igre) i tako je i navedeno u raspravi
(Poglavlje~\ref{Rasprava}).

Zaključak koji motivira sljedeći korak: nijedan prostorni filtar ne može
rekonstruirati informaciju koje u jednom ulaznom okviru nema — razlika među
njima je samo koju vrstu pogreške svaki radi. Tek akumulacija kroz vrijeme
donosi nove uzorke.

\section{Minimalni temporalni upscaler (TAAU)}

Drugi korak prvi put stvarno \emph{rekonstruira} razlučivost umjesto da je
interpolira: kamera svaki okvir pomiče projekciju za pod-pikselni Halton(2,3)
pomak, i taj niz pomaka jest dodatna informacija koju nijedan prostorni
filtar ne može izmisliti. Implementacija je namjerno minimalna — tri dijela:
reprojekcija, akumulacija i clamp. Dilatacija vektora, odbacivanje po dubini,
lockovi i Lanczos rekonstrukcija dolaze u punom upscaleru (\S\ref{fsr_full}).

\textbf{Reprojekcija.} Povijest (prethodni izlaz upscalera, u prikaznoj
razlučivosti) dohvaća se na \verb+uv + mv(uv)+, gdje vektor gibanja dolazi iz
G-buffera računat iz \emph{nejittriranih} clip-pozicija dok se rasterizacija
odvija s jittriranom projekcijom, i čita se isključivo točkastim uzorkovanjem
(filtrirani vektor na rubu objekta je prosjek dvaju nepovezanih gibanja i
pokazuje nikamo). Povijest se resamplira Catmull-Rom filtrom (9 bilinearnih
dohvata, standardna faktorizacija $4\times4$ jezgre s $a=-0.5$), što je
mjereno bolje od bilinearnog resamplinga za 0,75~dB u gibanju.

\textbf{Akumulacija.} Za svaki prikazni piksel uzima se $3\times3$ okolina
render piksela oko trenutne uzorkovane pozicije i teži se Gaussovom jezgrom
$\exp(-k d^2)$; rezultat se miješa s povijesti težinom koja raste do gornje
granice broja okvira. Širina jezgre je bio prvi stvarni nalaz razvoja: uska
jezgra ($k\approx6$--$8$) čini svaki pojedinačni okvir šumovitim, a taj šum
je upravo Halton uzorak koji se prosjekom pretvara u razlučivost — sužavanje
jezgre je vrijedilo $+0,5$~dB.

\textbf{Clamp.} Povijest izvan lokalnog raspona boja tekućeg okvira povlači
se natrag u njega, u granicama srednje vrijednosti $\pm\,\gamma$ standardnih
devijacija $3\times3$ okoline. Drugi stvarni nalaz razvoja: clamp mora biti u
\textbf{YCoCg} prostoru boje i mora biti \emph{clip} (povlačenje po pravcu
prema središtu kutije), ne nezavisan clamp triju RGB kanala — nezavisno
ograničavanje kanala pomiče ton (na svijetlom kamenu s tamnom fugom povuče
crvenu a ne plavu, i piksel ispadne magenta). Prelazak je vrijedio $+1,1$~dB
i uklonio vidljiv obrub.

\paragraph{Zašto TAAU u brzom gibanju gubi od bicubica.} Ovo je glavni
analitički nalaz koraka, ne bug. Interpolirajući filtar za resampliranje ima
nulu na Nyquistovoj frekvenciji kad pogodi pol-pikselnu fazu; okvir gibanja
koji nije cijeli broj piksela zato iz povijesti skida najfiniji detalj. Kako
se povijest svaki okvir vraća u akumulaciju s težinom $\approx(1-1/f)$, gubitak
po okviru postaje stacionaran:
\begin{equation}
G_{\text{stac}} \approx \frac{a}{1-(1-a)r}, \qquad a = \frac{1}{f},
\label{eq_nyquist}
\end{equation}
gdje je $a$ udio novog okvira u akumulaciji a $r$ prijenos filtra na
Nyquistu za tu fazu. Za $f=8$ i $r=0,5$ formula~\eqref{eq_nyquist} daje
$G_{\text{stac}}\approx0,22$ — više od dvije trećine najviših frekvencija
nestane, iako svaki pojedinačni resample gubi ``samo'' pola. Zato duga
povijest izgleda oštro dok kamera stoji i omekša čim se pomakne.

Posljedica za dizajn je \textbf{skraćivanje povijesti s gibanjem}: gornja
granica akumulacije pada eksponencijalno s brzinom gibanja u slici. Ablacija
tog mehanizma (Sponza, 60~fps orbita) pokazuje pomak od 2,5~dB između
isključenog i uključenog mehanizma — najveći pojedinačni pomak u cijeloj
ablaciji ovog koraka — po cijenu temporalne stabilnosti, koja se kreće u
suprotnom smjeru: fiksna granica akumulacije daje najstabilniju sliku u
tablici, i ujedno najlošiju vjernost, jer ta slika jednostavno ne prati
scenu.

\paragraph{Rezultati.} Sponza, 1280$\times$720~$\to$~1920$\times$1080
(Quality 1,5$\times$), referenca supersamplirana:

\begin{table}[hbtp]
\centering
\begin{tabular}{|l|c|c|c|}\hline
Kamera & Bicubic & TAAU & TAAU + RCAS \\\hline
mirna & 35,83 / 0,9441 & 39,36 / 0,9723 & \textbf{40,39 / 0,9819} \\\hline
120 fps orbita & 37,18 / 0,9596 & 37,50 / 0,9643 & \textbf{37,71 / 0,9690} \\\hline
60 fps orbita & \textbf{35,37} / 0,9420 & 35,13 / 0,9449 & 35,25 / \textbf{0,9505} \\\hline
30 fps orbita & \textbf{35,27} / 0,9392 & 34,36 / 0,9372 & 34,40 / \textbf{0,9421} \\\hline
\end{tabular}
\caption{TAAU naspram bicubica, po brzini kamere (PSNR dB / SSIM).}
\label{table_m4_rezultati}
\end{table}

Na mirnoj i sporo pomičnoj kameri TAAU+RCAS dobiva $+4,56$~dB odnosno
$+0,53$~dB nad najboljim prostornim filtrom — akumulacija radi ono zbog čega
postoji. SSIM je bolji na svakoj brzini, uključujući one na kojima PSNR
zaostaje: to je razlika u strukturi, ne u prosječnoj grešci — bicubic na
brzoj kameri ima puzajući aliasing na rubovima koji SSIM detektira, a PSNR
usrednji. Kriterij prihvaćanja iz plana rada (``slika stabilna, vidljiv
dobitak PSNR-a nad bicubic'') je ispunjen.

Zanimljivo je da se \textbf{RCAS kod TAAU-a ponaša suprotno} nego kod FSR1-a:
postoji pravi optimum oštrine ($\approx$1,0--1,5) iznad isključenog RCAS-a
($+0,13$~dB, $+0,0057$ SSIM na 60~fps; $+1,03$~dB na mirnoj kameri), jer
izoštravanje ovdje vraća Nyquist koji je resample povijesti dokazivo odnio, a
ne izmišlja detalj kojeg nema — suprotno FSR1-u, gdje EASU već izvlači sve
što jedan okvir sadrži.

\section{Puni upscaler (FSR3-stil)}\label{fsr_full}

Treći korak dodaje mehanizme koji minimalni TAAU čine upotrebljivim u
proizvodnom smislu: dilataciju vektora gibanja, odbacivanje povijesti po
dubini (disokluzija), lockove za tanke detalje, adaptivnu Lanczos
rekonstrukciju i reactive masku. Implementiran je kao tri compute prolaza:
dilatacija i depth-clip zajedno (dijele istu $3\times3$ okolinu dubine),
detekcija lockova, i glavni prolaz koji spaja Lanczos rekonstrukciju,
reprojekciju, upravljanje lockovima, rektifikaciju boje i akumulaciju.

\textbf{Dilatacija vektora.} Vektor gibanja vrijedi za površinu koju piksel
stvarno pokriva; na silueti render piksel pokriva dvije površine a sprema
samo jednu. Uzimanje vektora \emph{najbliže} dubine u $3\times3$ prozoru
pomiče svaki rubni piksel prema prednjoj (vizualno dominantnoj) površini.

\textbf{Depth clip.} G-buffer po pikselu piše i ovokvirnu i prethodnu
linearnu udaljenost iste površine (oboje besplatno uz reverse-Z projekciju s
beskonačnom dalekom ravninom, jer je \verb+clip.w = -viewZ+ egzaktno). Ako
vektor pokazuje na površinu koja je stvarno bila tamo prošli okvir, dubinski
spremnik na toj poziciji odgovara; ako ne, piksel je bio skriven i njegova
povijest je tuđa boja. Prag tolerancije je relativan i uvećan za lokalni
nagib dubine, da se pod gledan pod oštrim kutom ne proglasi lažno
disokludiranim.

\textbf{Lockovi.} Rektifikacija boje neselektivno povlači povijest prema
susjedstvu tekućeg okvira, što je kobno za tanke detalje (karika lanca,
stabljika, rub draperije) koji na neke Halton faze jednostavno ne padnu na
render uzorak: clamp tada briše akumulirani dokaz da detalj postoji i slika
treperi u ritmu jitter sekvence. Detekcija (po uzoru na FSR2) traži piksel
koji je strogi ekstrem luminancije duž barem jednog od četiri smjera kroz
$3\times3$ prozor — linija široka jedan piksel je ekstrem \emph{poprijeko}, a
ne \emph{uzduž}, što ga razlikuje od šuma (ekstrem u svakom smjeru) i od
ruba (ekstrem u nijednom). Lock živi ograničen broj okvira, putuje s
poviješću i dok traje širi clamp kutiju te joj briše presjek s tvrdim
min/max — jer je upravo taj presjek operacija koja briše detalj.

\textbf{Lanczos rekonstrukcija prilagođena gibanju.} Umjesto Gaussove jezgre
iz TAAU-a, glavni prolaz koristi Lanczos-2 preko $4\times4$ render uzoraka.
Širina jezgre vezana je na istu mjeru gibanja kao i duljina akumulacije:
\begin{verbatim}
stillness    = exp(-motionDecay * mvPixels);
maxWeight    = mix(1.0, maxWeightConst, stillness);
lanczosScale = mix(lanczosScaleMoving, lanczosScaleStill, stillness);
\end{verbatim}
Uska jezgra trguje oštrinu jednog okvira za njegov šum, što je bezopasno
samo dok ima dovoljno okvira da se usrednji. Izmjereni optimum na mirnoj
kameri je 2,0, u brzom gibanju 1,0.

\paragraph{Rezultati.} Ista scena i postavka kao TAAU:

\begin{table}[hbtp]
\centering
\begin{tabular}{|l|c|c|c|}\hline
Kamera & TAAU (prethodni korak) & FSR (ovaj korak) & $\Delta$ \\\hline
mirna & 39,36 / 0,9723 & \textbf{40,06 / 0,9828} & $+0,70$~dB \\\hline
120 fps & 37,50 / 0,9643 & 37,53 / 0,9636 & $+0,03$~dB \\\hline
60 fps & 35,13 / 0,9449 & \textbf{35,48 / 0,9460} & $+0,35$~dB \\\hline
30 fps & 34,36 / 0,9372 & \textbf{34,95 / 0,9388} & $+0,59$~dB \\\hline
\end{tabular}
\caption{Puni upscaler naspram TAAU, po brzini kamere.}
\label{table_m5_rezultati}
\end{table}

Najvažnija posljedica: na 60~fps orbiti puni upscaler \textbf{prvi put
nadmašuje bicubic i po PSNR-u i po SSIM-u} (35,48/0,9460 naspram
35,37/0,9420 iz \S 5.1) — minimalni TAAU je tamo još zaostajao. Temporalna
stabilnost se pri tom popravlja time što se \emph{približi} referenci, ne
time što raste: na 60~fps orbiti referenca ima stabilnost 34,00~dB, TAAU
36,94 ($+2,94$, previše stabilno — akumulacija prigušuje i legitimne
promjene), a puni upscaler 34,15 ($+0,15$, ponaša se temporalno kao ground
truth).

Ablacija po mehanizmu (60~fps orbita naspram mirne kamere) pokazuje da
\textbf{dilatacija donosi najveći pojedinačni dobitak, i to na mirnoj kameri}
($+2,50$~dB) — naizgled kontraintuitivno, jer se kamera ne miče, ali baš zato
je povijest najdulja i pogrešan vektor na silueti akumulira grešku kroz svih
osam okvira umjesto da je jednom napravi. Depth clip je mjerljivo koristan
samo u gibanju ($+0,13$~dB), a lockovi rade suprotno: $-0,06$~dB u gibanju,
$+0,51$~dB na miru.

Dva mehanizma iz nacrta dala su \textbf{negativan ili nulti rezultat, i oba
su izvještena kao takva}: reactive maska monotono šteti (Sponza nema pravu
prozirnost — jedine nekontinuirane površine su alfa-testirano lišće, čiji je
alfa test deterministički, pa je povijest koju maska baca bila točna), a
RCAS na izlazu punog upscalera sad šteti umjesto da pomaže, jer je
adaptivna Lanczos jezgra i lockovi već uklonili nedostatak (izgubljeni
Nyquist) koji je RCAS kod TAAU-a popravljao.

\textbf{Nedostaje} luma piramida iz izvornog nacrta (FSR2 je koristi za
auto-ekspoziciju) — ovaj pipeline tone-mapa \emph{prije} upscalera, pa je
ekspozicija već fiksirana i piramida ne bi promijenila ništa; izostavljena
je iz istog razloga i reverzibilni interni tonemap. Oba su namjerna
odstupanja od nacrta, ne propust, i raspravljena su u
Poglavlju~\ref{Rasprava}.


\chapter{Implementacija procjene gibanja}\label{ImplOpticalFlow}

Upscaler iz prethodnog poglavlja koristi vektore gibanja koje daje
rasterizacija — poznati su egzaktno za svaki trokut koji se crta. Oni,
međutim, ne opisuju sjene, refleksije, promjene sjenčanja, animirane
teksture niti sadržaj koji se ne renderira po trokutima (UI, čestice).
Procjena gibanja \emph{iz same slike} — optical flow — postoji upravo za taj
ostatak, i to je jezgra ``napredne metode računalnog vida'' u ovom radu.
Koristi ga isključivo generiranje okvira (Poglavlje~\ref{ImplFrameGen});
upscaler ga po zadanom ne koristi, da ne dodaje svoju cijenu svakom
mjerenju upscalera.

\section{Piramidalni block matching}

Implementacija procjenjuje gibanje između prikazanog i prethodnog
prikazanog okvira, za svaki blok od $8\times8$ texela luminancije (na pola
prikazne razlučivosti), po uzoru na \verb+FfxOpticalFlow+ iz FidelityFX SDK-a.
Vektori su, po konvenciji cijelog projekta, u UV prostoru i pokazuju iz
trenutnog okvira u prethodni — identično polju gibanja iz G-buffera, čime su
dva polja izravno usporediva i zamjenjiva.

Glavna petlja radi coarse-to-fine, 7 razina piramide, tri prolaza po razini:

\begin{enumerate}
\item \textbf{Pretraga} (\emph{search}) — za svaki nepreklapajući blok od
$8\times8$ texela, pretraga prozora radijusa $R$ oko predikcije iz grublje
razine, kriterij minimalna suma apsolutnih razlika (SAD) luminancije.
\item \textbf{Filtriranje} (\emph{filter}) — $3\times3$ vektorski medijan
(vektor s minimalnom sumom udaljenosti do preostalih osam) uklanja izolirane
ispade.
\item \textbf{Profinjenje} (\emph{upscale}) — prijelaz na finiju razinu bira
između kandidata (roditelj i tri susjeda) po SAD-u; ne izvodi se na zadnjoj
razini.
\end{enumerate}

Vektori se u ovoj implementaciji nose u UV prostoru \emph{pune slike} na
svakoj razini piramide, ne u pikselima razine, čime prijelaz između razina
postaje čist izbor kandidata bez množenja s dva — cijela klasa grešaka ``u
kojoj je razini ovaj vektor'' ne može nastati. Cijena je odsutnost sub-pixel
profinjenja: najsitniji razlučiv pomak je jedan texel luminancije na pola
razlučivosti, tj. dva prikazna piksela.

Kriterij podudaranja je SAD, a ne SSD (suma kvadrata razlika): kvadriranje bi
jednom specularnom pikselu s velikom promjenom dalo moć da poništi inače
savršeno podudaranje. SAD se normalizira u srednju apsolutnu razliku po
texelu, uz kaznu udaljenosti od predikcije (kazna glatkoće) u istim
jedinicama — bez toga bi pretraga unutar ravne plohe (gdje svaki kandidat
postiže isti SAD) birala slučajnog pobjednika umjesto grublje razine kao
prirodnog prior-a.

\section{Referenca i točnost}

Optical flow u pravom slučaju uporabe nema ground truth. Ali dok je jedino
gibanje gibanje kamere u statičnoj sceni (Sponza), vektori iz rasterizacije
su egzaktni i postaju besplatna referenca. Pogreška se mjeri kao endpoint
error (EPE) u prikaznim pikselima — euklidska udaljenost između procijenjenog
vektora i prosjeka referentnog polja preko površine bloka.

\begin{table}[hbtp]
\centering
\begin{tabular}{|l|c|c|c|c|c|}\hline
Kamera & EPE sred. & EPE medijan & EPE p95 & unutar 1 px & unutar 2 px \\\hline
mirna & 0,00 & 0,00 & 0,00 & 100,0\% & 100,0\% \\\hline
120 fps & 2,82 & 1,49 & 7,62 & 65,7\% & 89,0\% \\\hline
60 fps & 11,63 & 3,66 & 36,81 & 56,9\% & 80,2\% \\\hline
30 fps & 35,79 & 8,83 & 106,72 & 47,5\% & 70,4\% \\\hline
\end{tabular}
\caption{Točnost procjene gibanja po brzini kamere.}
\label{table_m6_brzina}
\end{table}

Za usporedbu, procjenitelj koji uvijek vrati nulti vektor daje srednju
pogrešku 23,0 / 46,1 / 92,3~px pri 120/60/30~fps — implementirani procjenitelj
dakle uklanja $\approx$87\% gibanja pri 120~fps, 75\% pri 60~fps i 61\% pri
30~fps. Degradacija s brzinom nije linearna jer block matching ima tvrd
doseg: radijus $R$ na razini $L$ pokriva $2^L\cdot R$ texela, a preko toga se
podudaranje ne pogoršava postupno nego se posve izgubi.

\section{Ablacije}

Tri mehanizma doprinose na različit način, mjereno na istoj orbiti (60~fps,
sve ostalo zadano):

\begin{table}[hbtp]
\centering
\begin{tabular}{|l|c|c|c|c|}\hline
Konfiguracija & EPE sred. & medijan & p95 & unutar 1 px \\\hline
sve zadano & 7,17 & 1,36 & 33,50 & 62,5\% \\\hline
bez kandidata iz prošlog okvira & 11,04 & 1,84 & 49,96 & 60,1\% \\\hline
bez medijan filtra & 7,95 & 2,39 & 34,64 & 56,5\% \\\hline
bez izbora kandidata pri profinjenju & 12,68 & 1,18 & 56,98 & 61,5\% \\\hline
\end{tabular}
\caption{Ablacija mehanizama procjene gibanja.}
\label{table_m6_ablacija}
\end{table}

\textbf{Temporalni kandidat} (polje prošlog okvira kao dodatni kandidat u
pretrazi) je najveći pojedinačni doprinos srednjoj pogrešci — gibanje je
vremenski koherentno, pa je prošlo polje bolji prior od bilo čega što
piramida sama može ponuditi, i ne degradira s veličinom pomaka kao piramida.
\textbf{Medijan filtar} se vidi u medijanu i udjelu točnih blokova, a jedva u
srednjoj vrijednosti — uklanja izolirane ispade, ne velike sistematske
promašaje. \textbf{Izbor kandidata pri profinjenju} je obrnut slučaj: bez
njega je medijan bolji, ali p95 naraste sa 33,5 na 57,0 — slijepo kopiranje
roditelja razmazuje granice grubljih blokova u finiju razinu.

Radijus pretrage je odabran kao kompromis točnosti i cijene: od 4 do 8
texela točnost se popravi 9\% uz 46\% skuplji modul, pa je 4 zadano.

\section{Detekcija promjene scene}

Nakon reza dva uzastopna okvira nemaju nikakve veze, a block matching bi
svejedno vratio ``najbolje'' podudaranje, koje je smeće. Detektor uspoređuje
histograme luminancije po devet sekcija slike, računa udaljenost totalne
varijacije po sekciji i iz devet brojeva izvodi jednu presudu. Mjereno na tri
orbite s umjetno umetnutim rezovima, statistika koja se uzima nad devet
sekcija odlučuje o pouzdanosti:

\begin{table}[hbtp]
\centering
\begin{tabular}{|l|c|c|c|c|}\hline
Statistika & min(rez) & max(gibanje) & detektirano & lažnih pozitiva \\\hline
maksimum & 0,568 & 0,463 & 15/15 & 11/105 \\\hline
\textbf{srednja} & 0,253 & 0,207 & 15/15 & \textbf{0/105} \\\hline
medijan & 0,232 & 0,179 & 14/15 & 0/105 \\\hline
\end{tabular}
\caption{Statistika detekcije reza scene, mjerna orbita.}
\label{table_m6_rez}
\end{table}

Maksimum ima najširu marginu i ipak je najgori izbor: brzi zaokret kamere
zamijeni sadržaj nekoliko sekcija na vodećem rubu, što se od stvarnog reza ne
razlikuje ni po čemu drugom, i proizvodi lažne pozive. Srednja vrijednost
hvata sve rezove uz nula lažnih poziva na dvije od tri testirane orbite i
zadana je statistika.

\section{Cijena i granice}

Modul košta $\approx$0,56~ms na 1080p (720p~$\to$~1080p upscaling) odnosno
$\approx$21\% cijene okvira; najfinija razina piramide je 41\% te cijene, šest
grubljih razina zajedno manje od trećine. Poznata ograničenja, izmjerena i
prihvaćena kao svojstvo metode: periodične teksture (fino tkanje, prugasti
stupovi) daju bloku mnogo jednako dobrih podudaranja pa je procjena tamo
šumovita (aperture problem u Fleet--Weiss formulaciji); ravne plohe bez
strukture nemaju jedinstven minimum i procjena pada na prior; a blok koji
prekriva siluetu s dva gibanja mora izabrati jedno, što je granica rezolucije
mreže bloka, ne greška algoritma.

Teorijska podloga je Fleet i Weiss, \emph{Optical Flow Estimation}
\cite{FleetWeiss}: pretpostavka konstantnosti svjetline, aperture problem i
loša postavljenost problema unutar homogenih područja — tri pojave koje
ablacije u ovom poglavlju izravno mjere.


\chapter{Implementacija generiranja okvira}\label{ImplFrameGen}

Zadnji dio klasičnog cjevovoda iz dva uzastopna upscalana okvira, vektora
gibanja iz G-buffera i polja optical flowa (Poglavlje~\ref{ImplOpticalFlow})
gradi okvir na trenutku $t-0,5$ — onaj koji bi se prikazao između njih — i
stavlja ga na ekran u pravom trenutku. Ovo poglavlje pokriva izgradnju
međuokvira (Modul D iz plana) i njegov prikaz: frame pacing, latenciju i
kompoziciju korisničkog sučelja.

\section{Izgradnja međuokvira}

Konvencija je ista kao u cijelom projektu: vektori pokazuju iz trenutnog
okvira u prethodni. Površina koja se između $t-1$ i $t$ pomaknula za vektor
$\vec{v}$ bila je na pola puta u $t-0,5$, pa se međuokvir iz trenutnog čita
na $uv - 0,5\vec{v}$, a iz prethodnog na $uv + 0,5\vec{v}$.

Postupak ide u devet compute prolaza:

\begin{enumerate}
\item \textbf{Postavljanje} — pražnjenje ciljnih slika (prioritet 0).
\item \textbf{Dubina međuokvira} — scatter dubine na $uv+0,5\vec v$, atomski minimum.
\item \textbf{Polje game vektora} — scatter vektora iz G-buffera s prioritetom.
\item \textbf{Piramida polja} — popunjavanje rupa u log($n$) koracima (zajednička za oba polja).
\item \textbf{Polje optical flow vektora} — isti scatter, iz procjene gibanja.
\item \textbf{Maske disokluzije} — dvije maske, međuokvir$\leftrightarrow$prethodni i međuokvir$\leftrightarrow$trenutni.
\item \textbf{Interpolacija} — dva warpa po polju, izbor između polja po slaganju boja; alfa kanal = pokrivenost.
\item \textbf{Piramida za inpainting} (M8) — piramida slike množene pokrivenošću.
\item \textbf{Inpainting} (M8) — popunjavanje preostalih rupa iz piramide.
\end{enumerate}

\textbf{Zašto scatter i atomici.} Polje vektora međuokvira ne postoji u
nijednom ulazu — svaki izvorni piksel piše tamo gdje je njegova površina u
$t-0,5$, i više piksela može pasti na isti texel odredišta. Pobjednik se bira
bez drugog prolaza kroz kandidate pomoću \verb+imageAtomicMax+ nad
32-bitnom riječi čiji gornjih 16 bita nosi prioritet a donjih 16 komponentu
vektora kao \emph{half float} — atomik uspoređuje prvo prioritet, vrijednost
tek kao razrješenje izjednačenja. Prioritet game polja kodira (od
najznačajnijeg bita): je li površina najbliža na tom texelu, dubinu
(bliže~=~više) i slaganje boje na kraju vektora; prioritet polja toka kodira
umjesto dubine sigurnost podudaranja bloka i magnitudu pomaka.

\textbf{Popunjavanje rupa} u piramidi zadržava, pri redukciji, dijete s
\emph{najnižim} prioritetom (pozadinu): piksel u disokluziji pripada onome
što je bilo \emph{iza} objekta koji se odmaknuo, i to opisuje vektor
pozadine, ne vektor prednje površine.

\textbf{Interpolacija} radi dva izbora: unutar polja (koji od dva warpa
vrijedi — određuju maske disokluzije) i između polja (kojem vektoru
vjerovati — određuje slaganje boja na krajevima warpa, s težinom
$\exp(-k\cdot\text{udaljenost boja})$; točan vektor dovede istu površinu na
samu sebe i oba warpa se slažu, pogrešan se sam prijavi razlikom).

\section{Rezultati}

Sve mjereno na Sponzi, referenca je pravi render trenutka $t-0,5\Delta t$
(supersampliran, bez jittera), a stupac ``blend'' je naivni 50/50 blend
istog para ulaznih okvira naspram iste reference — cijena nečinjenja i
jedina poštena usporedba (ista razlučivost, isti upscaler, isti trenutak).

\begin{table}[hbtp]
\centering
\begin{tabular}{|l|c|c|c|}\hline
Konfiguracija & PSNR (min) & PSNR blend & dobitak \\\hline
1080p Quality & 34,94 (27,56) & 26,11 & $+8,83$~dB \\\hline
1080p native & 36,43 (27,92) & 26,14 & $+10,29$~dB \\\hline
1080p Performance & 33,52 (27,51) & 26,16 & $+7,36$~dB \\\hline
\end{tabular}
\caption{Interpolirani okvir naspram 50/50 blenda, po razlučivosti render prolaza.}
\label{table_m7_razlucivost}
\end{table}

Dobitak nad blendom je 7--10~dB na svakoj razlučivosti; kvaliteta međuokvira
prati kvalitetu ulaza (native vodi, Performance zaostaje). Po brzini kamere,
mirna kamera je provjera ispravnosti (blend je egzaktan, modul ga ne kvari
do stotinke), a srednji dobitak je gotovo neovisan o brzini gibanja
($+8$ do $+9$~dB), ali \textbf{najgori okvir} nije: pri 30~fps pada na razinu
prosjeka blenda, pri 20~fps ispod nje — pretpostavka linearnog gibanja i sve
šire disokluzije tamo prestaju vrijediti.

Ablacija na dvije scene (Sponza, gdje su game vektori egzaktni, i procedu\-ralna
scena s pomičnim objektima) pokazuje da se \textbf{dva polja nadopunjuju}:
samo game vektori daju 34,80/33,68~dB (Sponza/proceduralna), samo optical
flow 33,74/33,53~dB, a oba zajedno 34,94/34,07~dB — nijedno polje samo ne
doseže kombinaciju. Maske disokluzije donose malen ali mjerljiv dobitak
($+0,09$~dB na Sponzi), dok dilatacija vektora i izbor djeteta u piramidi
nemaju mjerljiv učinak na ovim scenama (rupe u polju su preuske da bi izbor
susjeda promijenio vektor) — što je i samo rezultat: mehanizmi postoje kao u
FSR3, na ovim scenama nemaju što popraviti.

\section{Prikaz: frame pacing, latencija, UI}

Interpolirani okvir mora stići na ekran u pravom trenutku, ne odmah kad je
gotov. Tri su pitanja, sva o vremenu i redoslijedu, ne o slici: \emph{kada}
se prikazuje koji okvir (pacing), \emph{koliko kasni} ono što igrač vidi za
onim što je napravio (latencija), i \emph{gdje u lancu} ulazi korisničko
sučelje.

\textbf{Arhitektura.} Render nit radi na skrivenom $1\times1$ prozoru;
zaseban nit prikaza, s vlastitim GL kontekstom koji dijeli teksture, brine se
o trenutku prikaza. Oba okvira para (generirani i pravi) crtaju se u
back buffer odmah pri predaji (\emph{staging}), pa u trenutku pravog prikaza
preostaje samo swap. Prva verzija (blit neposredno prije swapa) davala je
okvire koji su i dalje stizali u parovima, jer dva konteksta na istom GPU-u
ne mogu preskočiti red poslova — staging to zaobilazi predajom oba bliteta
odmah dok je GPU prazan.

\textbf{Rezultati pacinga.} Sponza, sintetsko opterećenje (dodatni G-buffer
prolazi po okviru simuliraju cijenu igre):

\begin{table}[hbtp]
\centering
\begin{tabular}{|l|c|c|c|c|}\hline
Opterećenje & bez FG & FG ravnomjerno & ubrzanje \\\hline
$\times$0 & 313,0 & 308,3 & 0,98$\times$ \\\hline
$\times$4 & 155,3 & \textbf{211,0} & 1,36$\times$ \\\hline
$\times$12 & 90,7 & \textbf{135,3} & 1,49$\times$ \\\hline
$\times$24 & 57,7 & \textbf{89,9} & 1,56$\times$ \\\hline
\end{tabular}
\caption{Prikazani FPS naspram sintetskog opterećenja.}
\label{table_m8_fps}
\end{table}

Ispod praga isplativosti generiranje ne donosi ništa ($\times$0: FPS ostaje
na mjestu, jer je cijena generiranja bliska cijeni renderiranja lagane
scene), a iznad njega ubrzanje raste na 1,5--1,6$\times$ — nikad na
2$\times$, jer generiranje nije besplatno i uvijek renderirani FPS uz njega
padne. \textbf{Ravnomjernost} pacinga je mjerena standardnom devijacijom
razmaka između prikazanih okvira: bez pacinga 7,03~ms (bimodalno, okviri
stižu u parovima), s ravnomjernim pacingom 1,00~ms — manje nego bez
generiranja uopće (1,38~ms).

\textbf{Latencija} raste za $\approx$jedan render period naspram okvira bez
generiranja (npr. $+11$~ms uz $\times$12, $+16$~ms uz $\times$24) — to je
cijena tehnologije, mjerena i izvještena, ne pretpostavljena.

\textbf{UI kompozicija.} HUD komponiran \emph{nakon} generiranja (na čistoj
slici bez teksta) naspram upečenog \emph{prije} njega (kao dio scene koju
warp pomiče): razlika je 32~dB u pravokutniku HUD-a i 4,4~dB na cijelom
okviru — upečeni HUD gubi više kvalitete nego sve ablacije generiranja
zajedno, jer warp tretira tekst kao dio scene i udvostruči ga.

Cijena čitavog generiranja okvira (uz optical flow) na 1080p Quality je
$\approx$2,2~ms, što je usporedivo s cijenom jednog pravog okvira ove
testne scene — otud prag isplativosti u tablici gore, i otud napomena da je
svaka scena teža od testne iznad tog praga.


\chapter{Naučena komponenta}\label{NaucenaKomponenta}

Opis teme traži i naučenu komponentu, a FSR3 je u cijelosti klasičan
računalni vid — naučeni komercijalni pristupi (DLSS, FSR4 ``Redstone'')
žive na hardveru s matričnim jedinicama koje razvojni GPU ovog rada nema
(Poglavlje~\ref{PregledRjesenja}). Ova komponenta zato ne zamjenjuje jezgru
cjevovoda, nego mjeri jedno izolirano pitanje: \textbf{može li mala mreža,
uz iste ulaze, bolje od ručne heuristike odlučiti kako spojiti kandidate
koje klasični lanac (Poglavlje~\ref{ImplFrameGen}) već proizvodi — i koliko
to košta na GPU-u bez podrške za tenzore.}

\section{Što mreža radi}

Mreža \textbf{ne crta piksele}. Za svaki piksel predviđa softmax težine nad
sedam kandidata boje koje klasični dio generiranja okvira već ima: izlaz
heuristike, warp igrinim vektorima iz $t-1$ i iz $t$, warp optical flowom
iz $t-1$ i iz $t$, te trenutni i prethodni okvir bez warpa. Izlaz je njihov
ponderirani zbroj. Tri razloga za takav oblik:

\begin{enumerate}
\item \textbf{Čista usporedba.} Ulazi i kandidati su identični heuristici;
naučeno je samo pravilo miješanja, pa je razlika u PSNR-u razlika pravila,
ne dodatne informacije.
\item \textbf{Ne može halucinirati.} Izlaz je konveksna kombinacija
postojećih boja, pa mreža ne treba kapacitet za sintezu detalja — staje u
nekoliko tisuća parametara i u compute shader.
\item \textbf{Počinje od heuristike.} Pristranost zadnjeg sloja je
inicijalizirana tako da kandidat 0 (heuristika) dobije $\approx$90\%
težine; svaki odmak od klasičnog rezultata učenje mora zaraditi na
podacima.
\end{enumerate}

\section{Arhitektura}

Mala U-Net-lite mreža (odabrani model, oznaka c8-c16, ima 7199 parametara)
radi u compute shaderima na 20 značajki po pikselu (valjanost i pokrivenost
polja, slaganje boja među kandidatima, udaljenosti heuristike od warpova,
razine piramide, duljine vektora, luminancija). Arhitektura ima enkoder od
tri razine (pooling $2\times2$ + konvolucija $3\times3$) i dekoder s
preskočnim vezama, završava $1\times1$ slojem i softmaxom. Kanali žive četiri
po RGBA16F sloju teksturnog polja, pa je $3\times3$ konvolucija $16\to16$
kanala oblik posla ($4\times9\times4$ umnožaka $\text{mat4}\cdot\text{vec4}$
po pikselu) u kojem su i GPU-ovi bez matričnih jedinica još razumno dobri.

\section{Skup podataka i učenje}

Skup podataka snima vlastiti renderer u lockstep načinu: svaki mjereni okvir
daje patcheve $96\times96$ značajki i kandidata, uz \textbf{stvarno
renderirani među-okvir} kao oznaku (isti postupak kao referenca u
Poglavlju~\ref{ImplFrameGen}). Podjela na skup za učenje i validaciju je po
\emph{putanji kamere}, ne po okviru — uzastopni okviri iste putanje su
gotovo ista slika, pa bi podjela po okviru ocjenjivala mrežu na podacima
koje je efektivno već vidjela. Mjerni pogledi korišteni za rezultate u
ovom radu nikad se ne snimaju u skup za učenje.

Trening koristi vlastiti trener napisan u C++20 s OpenMP-om, na CPU-u, a ne
PyTorch na vanjskom GPU-u: model je dovoljno malen (nekoliko tisuća
parametara) da su unaprijedni i povratni prolaz nekoliko stotina linija
petlji koje prevoditelj vektorizira, a razvojni GPU nema ROCm podršku
potrebnu za trening na njemu. Posjedovanje obje strane (trener i shader
inferencije) omogućuje i međusobnu provjeru: snimka napravljena s
učitanom mrežom nosi u svakom patchu i izlaz shadera, a trener izračuna
vlastiti unaprijedni prolaz i uspoređuje ga s njim — za odabrani model
srednja apsolutna razlika je $7,1\times10^{-5}$, najveća $4,7\times10^{-4}$
(zaokruživanje half-floata). Bez ove provjere greška u redoslijedu stupaca
matrice ili u poravnanju poolinga davala bi mrežu koja naizgled radi, a
tiho gubi kvalitetu.

\section{Rezultati}

\begin{table}[hbtp]
\centering
\begin{tabular}{|l|c|c|c|}\hline
Pogled & Heuristika (PSNR/SSIM) & Mreža (PSNR/SSIM) & Dobitak \\\hline
1080p Quality & 35,00 / 0,9443 & 35,10 / 0,9458 & $+0,10$~dB \\\hline
1080p native & 36,50 / 0,9628 & 36,69 / 0,9650 & $+0,19$~dB \\\hline
proceduralna scena & 34,06 / 0,9708 & 34,26 / 0,9722 & $+0,20$~dB \\\hline
20 fps & 30,95 / 0,9122 & 31,11 / 0,9181 & $+0,16$~dB \\\hline
\end{tabular}
\caption{Naučena mješavina naspram heuristike, odabrani mjerni pogledi.}
\label{table_m9_rezultati}
\end{table}

SSIM raste na svih osam mjerenih pogleda, PSNR na sedam od osam
($+0,03$ do $+0,20$~dB), stabilno preko tri sjemena treniranja. Dobitak
\textbf{nije sustavan na najgorem pojedinačnom okviru} (mreža uči prosjek;
rijetki okviri s brzim zamahom kamere su u skupu za učenje zastupljeni
koliko su i rijetki) i \textbf{ne prenosi se na neviđenu scenu}: model
treniran samo na jednoj sceni daje na drugoj, njemu neviđenoj sceni,
praktički identičan rezultat heuristici.

\section{Cijena}

Naučeni dio na 1080p Quality (Sponza) košta 2,7~ms po okviru GPU vremena
(generiranje okvira ukupno 1,84~$\to$~4,58~ms). Pri sintetskom opterećenju
$\times$12 (Poglavlje~\ref{ImplFrameGen}, odjeljak o pacingu) prikazani FPS
pada s 147 na 120. \textbf{Na razvojnom hardveru se ne isplati}: heuristika
je 0,10~dB lošija za 40\% cijene generiranja okvira. To je izmjeren, ne
pretpostavljen odgovor na pitanje teme: gdje je granica malog naučenog
modela bez tenzorskih jedinica.

\section{Demonstracija}

Slika~\ref{fig_m9_20fps} pokazuje oba ekstrema u istom okviru (20~fps,
najveća disokluzija): u jednom izrezu mreža uklanja prugavi artefakt na
rubu (15,21~$\to$~38,15~dB), u drugom zamuti sitnu granu lišća
(25,11~$\to$~21,03~dB). Slika~\ref{fig_m9_1080p} pokazuje isti obrazac u
uobičajenom radnom režimu (1080p Quality), a Slika~\ref{fig_m9_proc} pokazuje
neviđenu scenu, gdje je razlika zanemariva — vizualna potvrda nalaza iz
prethodnog odjeljka.

\begin{figure}[h!]
\centering
\includegraphics[width=\textwidth]{../../captures/ml/gallery/sponza-20fps-frame18.png}
\caption{20~fps, okvir 18: najveći dobitak i najveći gubitak mreže u istom
okviru. Stupci: referenca, heuristika, naučena mješavina, prikaz težina,
pogreška heuristike i mreže ($\times$6).}
\label{fig_m9_20fps}
\end{figure}

\begin{figure}[h!]
\centering
\includegraphics[width=\textwidth]{../../captures/ml/gallery/sponza-1080p-q-frame57.png}
\caption{1080p Quality, okvir 57: tipičan radni režim.}
\label{fig_m9_1080p}
\end{figure}

\begin{figure}[h!]
\centering
\includegraphics[width=\textwidth]{../../captures/ml/gallery/proceduralna-frame67.png}
\caption{Proceduralna scena (neviđena pri treniranju), okvir 67: razlika
zanemariva.}
\label{fig_m9_proc}
\end{figure}

\begin{figure}[h!]
\centering
\includegraphics[width=0.4\textwidth]{../../captures/ml/gallery/legend.png}
\caption{Legenda boja prikaza težina iz slika iznad.}
\label{fig_m9_legend}
\end{figure}


\chapter{Metodologija mjerenja}\label{Metodologija}

\section{Ground truth}

Kvaliteta u ovom radu je brojka jer postoji referenca, a referenca postoji
zato što vlastiti renderer to dopušta. Za upscaler, referenca je puni okvir
u prikaznoj razlučivosti, bez jittera, u trenutku $t$. Za generiranje okvira,
referenca je scena i kamera pomaknuti na $t-0,5\Delta t$ (bez upisivanja u
povijest instanci, da idući pravi okvir ne dobije prekratke vektore
gibanja), renderirani bez jittera i vraćeni na $t$. Oba referentna prolaza
renderiraju se supersamplirano ($2\times2$) i usrednjena u linearnom
prostoru boje prije tone mappinga — jednostruka referenca u prikaznoj
razlučivosti bi bila aliasirana, a bodovanje naspram aliasirane referencije
sustavno nagrađuje zamućenje (Poglavlje~\ref{TeorijskaPodloga}).

Ispravnost među-okvira provjerava se neovisno o glavnom mjernom putu:
pomoćna skripta provjerava da je referentni među-okvir podjednako udaljen
(PSNR) od oba susjedna pravog okvira i bliži svakom od njih nego što su
susjedi jedan drugome. Na testnoj sekvenci asimetrija je 0,01~dB — među-okvir
leži na sredini unutar mjerne pogreške. Ovom provjerom dobiva se i donja
granica koju generiranje okvira mora nadmašiti da bi imalo smisla: PSNR
naivnog 50/50 blenda dva susjedna pravog okvira naspram istog referentnog
među-okvira (28,85~dB na testnoj putanji) — donja granica jer i blend i
generiranje okvira koriste isti ulazni par, a razlika je što blend ne zna
ništa o gibanju.

\section{Skriptirana kamera i determinizam}

Svako mjerenje koristi skriptiranu putanju kamere i fiksni vremenski korak
($\Delta t$), umjesto stvarnog proteklog vremena — bez toga dva pokretanja
ne bi proizvela istu sekvencu okvira, a referenca koja se ne može ponoviti
nije referenca. Način rada s fiksnim brojem okvira ne prihvaća unos s
tipkovnice ni miša: jedan slučajan pritisak tijekom mjerenja (npr.
razmaknica, koja pauzira scenu) tiho bi promijenio ono što se mjeri bez
greške u izvođenju programa. Mjerna skripta zato provjerava da se vrijeme
scene stvarno pomiče između uzastopnih okvira mjerenja i prekida mjerenje
ako ne — jedina iznimka je mjerenje s namjerno mirnom kamerom, gdje se
vrijeme scene pomiče za korak ispod razlučivosti izvještaja.

\section{Mjerenje performansi}

Vrijeme na GPU-u mjeri se \emph{po passu}, GPU timestamp upitima (ne CPU
mjerenjem obrađenim naknadno), jer je upravo raspodjela troška po passu ono
što rad čini inženjerskim rezultatom, ne demonstracijom. Prvih nekoliko
okvira svakog mjerenja (\emph{warmup}) se odbacuje iz izvještaja: brojaču
vremena treba nekoliko okvira da počne vraćati broj (prsten upita), a prvi
pravi okviri uključuju prevođenje shadera i prve alokacije tekstura, pa
mogu biti red veličine skuplji od stabilnog stanja. Prolazi koji postoje
samo radi mjerenja (npr. renderiranje dodatne referentne slike uz
\verb+--validate-upscale+) mjere se, ali se eksplicitno isključuju iz
ukupne cijene okvira koja se izvještava — inače bi mjerenje samo sebe
naplaćivalo.

\section{Mjerenje kvalitete: PSNR i SSIM na GPU-u}

PSNR i SSIM računaju se izravno na GPU-u, u compute shaderu, u dvostupanjskoj
redukciji: prvi prolaz akumulira djelomične sume po radnim skupinama
$16\times16$ u storage buffer, a CPU tek zbraja te djelomične sume — izravna
atomska akumulacija u \verb+float+ preko cijele slike bi izgubila
preciznost (velik broj malih dodataka istoj varijabli). SSIM se računa po
formulaciji Wanga i sur. \cite{WangSSIM}: Gaussov prozor $11\times11$
($\sigma=1,5$) na kanalu luminancije, konstante $C_1=0,0001$,
$C_2=0,0009$, u dva separabilna prolaza (vodoravni i uspravni).

Temporalna stabilnost (koliko slika treperi kroz vrijeme, što PSNR/SSIM
jednog okvira ne mjeri) izvještava se \textbf{relativno naspram iste mjere
nad referentnim nizom}, ne naspram nule: promjene sjenčanja, disokluzije i
resampliranje same referencije prežive reprojekciju kao i izlaz upscalera,
pa referenca nije savršeno stabilna nula. Iznad $\approx$50~dB ova metrika
prestaje razlikovati bilo što ($\text{RMS}$ razlika $<0,002$) i u tom
području se ne čita.

\section{Alat za regeneraciju tablica}

Svih 283 konfiguracija koje se pojavljuju u ovom radu regenerira se jednom
naredbom (\verb+scripts/run_metrics.py+): skripta redom pokreće svaku
konfiguraciju, prikuplja CSV po okviru i agregira sve rezultate koji se
trenutno nalaze na disku u jednu zbirnu tablicu — bez obzira je li ih
pokrenula upravo ta invokacija ili prijašnja. To znači da se cijeli mjerni
paket, na bilo kojem računalu, može reproducirati bez ručnog prepisivanja
brojki, što je i korišteno za provjeru na dva različita GPU-a u
odjeljku~\ref{ProvjeraHardver}.


\chapter{Rezultati}\label{Rezultati}

\section{Provjera na drugom hardveru}\label{ProvjeraHardver}

Cijeli mjerni paket (283 konfiguracije, \verb+scripts/run_metrics.py+) je
pokrenut na dva potpuno različita računala, s istim izvornim kodom (isti
\emph{commit}), da se provjeri da rezultati ne ovise o jednom proizvođaču
GPU-a niti o slučajnosti jedne mjerne postavke.

\begin{table}[hbtp]
\centering
\begin{tabular}{|l|l|l|}\hline
 & Računalo A & Računalo B \\\hline
GPU & NVIDIA GeForce RTX 5070 & AMD Radeon RX 7800 XT (Sapphire Nitro+) \\\hline
CPU & AMD Ryzen 7 7800X3D & AMD Ryzen 7 5700X3D \\\hline
Driver & NVIDIA 616.92 (32.0.16.1692) & najnoviji AMD Adrenalin \\\hline
OS & Windows & Windows \\\hline
\end{tabular}
\caption{Dvije mjerne postavke korištene za provjeru.}
\label{table_hw_setup}
\end{table}

\textbf{Ispravnost.} Sve 283 konfiguracije podudaraju se po imenu između dva
skupa rezultata, pa je usporedba red po red. Metrike kvalitete (PSNR, SSIM,
endpoint error procjene gibanja) gotovo se ne razlikuju: srednja apsolutna
razlika je reda $10^{-3}$--$10^{-1}$ dB odnosno piksela na svim mjerenim
stupcima. Najveće pojedinačne razlike (do $\approx$0,7--0,9 dB) padaju
isključivo na konfiguracije koje su već u poglavljima~\ref{ImplUpscaler} i
\ref{ImplOpticalFlow} obilježene kao mjerno osjetljive same po sebi — mirna
kamera (PSNR iznad $\approx$40 dB, gdje se metrika više ne razlikuje
smisleno, vidi napomenu u odjeljku o temporalnoj stabilnosti) i ekstremna
ablacijska postavka \verb+RCAS sharpness = 0+. Nigdje se ne javlja
sustavno odstupanje karakteristično za, primjerice, grešku u redoslijedu
operacija ili u pretpostavci o preciznosti pokretnog zareza specifičnoj za
jednog proizvođača. Zaključak: implementacija se ponaša determinističko i
prenosivo preko granice AMD/NVIDIA, što nije samo po sebi razumljivo — GLSL
prevoditelji, redoslijed zbrajanja u sumama i ponašanje teksturnih filtara
razlikuju se između proizvođača.

\textbf{Performanse.} Ukupni omjer GPU vremena (RX 7800 XT / RTX 5070) je
blizu 1 (medijan 1,07, srednja vrijednost 1,02) — dva su GPU-a na ovom
opterećenju vrlo izjednačena — ali raspodjela po modulu pokazuje dosljedan
obrazac, ne slučajan šum:

\begin{table}[hbtp]
\centering
\begin{tabular}{|l|c|l|}\hline
Modul & Medijan omjera & Čitanje \\\hline
Optical flow (M6) & 0,85--0,88 & RX 7800 XT $\approx$15\% brži \\\hline
Naučena mješavina (M9) & 0,88--0,93 & RX 7800 XT $\approx$10\% brži, više na većim modelima \\\hline
TAAU, jezgra generiranja okvira (M4, M7) & 1,00--1,08 & praktički izjednačeno \\\hline
Puni upscaler FSR (M5) & 1,08--1,14 & RX 7800 XT $\approx$10\% sporiji \\\hline
\end{tabular}
\caption{Omjer GPU vremena po modulu, RX 7800 XT naspram RTX 5070.}
\label{table_hw_speed}
\end{table}

Obrazac je čitljiv arhitekturno: AMD je brži na prolazima s mnogo aritmetike
po texelu i malo teksturnih dohvata (block matching u optical flowu,
$3\times3$/$4\times4$ konvolucije naučene mješavine), a NVIDIA na
akumulacijskom prolazu punog upscalera, koji je dominantno vezan na filtrirane
teksturne dohvate (Catmull-Rom i Lanczos uzorkovanje povijesti). Nijedan GPU
nije jednoznačno brži na cijelom cjevovodu, što je i razlog zašto se
apsolutne brojke u ovom radu izvještavaju po modulu, a ne kao jedna
sažeta brojka za ``FPS''.

\section{Sinteza po modulu}

Detaljni rezultati i ablacije po modulu, s tablicama i analizom, dani su uz
odgovarajuću implementaciju: prostorni baseline, TAAU i puni upscaler u
Poglavlju~\ref{ImplUpscaler}, procjena gibanja u Poglavlju~\ref{ImplOpticalFlow},
generiranje okvira i prikaz u Poglavlju~\ref{ImplFrameGen}, naučena
komponenta u Poglavlju~\ref{NaucenaKomponenta}. Ovaj odjeljak sažima
headline brojke na jednom mjestu, redoslijedom kojim se cjevovod nadograđuje.

\begin{table}[hbtp]
\centering
\small
\begin{tabular}{|l|l|c|}\hline
Korak & Mjera dobitka & Vrijednost \\\hline
M3 (FSR1 EASU) naspram nearest & PSNR & $+3,27$~dB \\\hline
M4 (TAAU+RCAS) naspram M3 bicubic, mirna kamera & PSNR & $+4,56$~dB \\\hline
M5 (puni upscaler) naspram M4, 60~fps orbita & PSNR & $+0,35$~dB \\\hline
M6 optical flow naspram nulte procjene, 120~fps & udio uklonjenog gibanja & $\approx$87\% \\\hline
M7 generiranje okvira naspram 50/50 blenda, 1080p Q. & PSNR & $+8,83$~dB \\\hline
M8 pacing naspram trenutnog prikaza, opterećenje $\times$12 & std. razmaka & $7,03\to1,00$~ms \\\hline
M9 naučena mješavina naspram heuristike & SSIM & $+$ na 8/8 pogleda \\\hline
\end{tabular}
\caption{Headline dobitak svakog koraka cjevovoda naspram prethodnog.}
\label{table_ch10_sinteza}
\end{table}

Svaki korak nadmašuje prethodni na istoj sceni i istoj metrici — kriterij koji
je u planu rada postavljen kao uvjet prihvaćanja za svaki milestone. Iznimka
koja se mora navesti otvoreno: M9 pomiče kvalitetu (SSIM raste na svih osam
mjerenih pogleda), ali ne pomiče omjer kvalitete i cijene na razvojnom
hardveru — rasprava o tome je u Poglavlju~\ref{Rasprava}.

\section{Latencija}

Generiranje okvira nužno podiže latenciju: pravi okvir $t$ mora se prikazati
istovremeno s (ili nakon) generiranim okvirom koji leži između $t-1$ i $t$,
pa ravnomjeran prikaz zahtijeva čekanje reda veličine pola render perioda.
Izmjereno (Poglavlje~\ref{ImplFrameGen}), latencija do pravog okvira raste
za $+11$ do $+16$~ms naspram prikaza bez generiranja, ovisno o opterećenju —
približno jedan render period, kako i analiza mehanizma predviđa. Ovo je
prikazano kao cijena tehnologije, ne sakriveno u prosjeku FPS-a.

\section{Provjera na dva GPU-a proizvođača}

Vidi odjeljak~\ref{ProvjeraHardver}: kvaliteta se poklapa preko AMD/NVIDIA
granice na $10^{-3}$--$10^{-1}$~dB, a performanse po modulu pokazuju
dosljedan arhitekturni obrazac (AMD brži na aritmetički-intenzivnim
prolazima procjene gibanja i naučene mješavine, NVIDIA brži na
teksturno-intenzivnom akumulacijskom prolazu upscalera), ne slučajan šum.

\section{Artefakti}

Vizualne usporedbe (uvećani izresci kroz sve upscalere, HSV prikaz procjene
gibanja naspram egzaktnog polja, usporedba interpoliranog okvira i 50/50
blenda, usporedba komponiranog i upečenog korisničkog sučelja) generiraju
zasebne skripte iz \verb+scripts/+ direktorija i nisu uključene ovdje kao
zaseban skup slika, osim demonstracije naučene komponente
(Slike~\ref{fig_m9_20fps}--\ref{fig_m9_legend}) koja je izravno relevantna za
zaključak Poglavlja~\ref{NaucenaKomponenta}.


\chapter{Rasprava i ograničenja}\label{Rasprava}

\section{Odstupanja od izvornog plana}

\textbf{Metrike kvalitete.} Izvorni plan predviđao je PSNR, SSIM, FLIP i
LPIPS. Implementirani su samo PSNR i SSIM, oba na GPU-u. FLIP i LPIPS
zahtijevaju predtrenirane perceptivne modele odnosno CIE-prostor razlike
percepcije koji nadilaze GPU compute shadere ovog projekta; PSNR/SSIM su se
u praksi pokazali dovoljno osjetljivi da razlikuju sve ablacije izvještene u
poglavljima~\ref{ImplUpscaler}--\ref{NaucenaKomponenta}, uz poznatu manu
(PSNR nagrađuje zamućenje na aliasiranom sadržaju, djelomično riješeno
supersampliranom referencom — Poglavlje~\ref{Metodologija}) koja je svugdje
gdje je relevantna i eksplicitno navedena uz brojku.

\textbf{Usporedba s Lucas--Kanade metodom.} Plan je predviđao opcijski
dodatak: implementaciju Lucas--Kanade gradijentne metode za procjenu gibanja
i usporedbu točnosti/brzine s implementiranim block matchingom. Ovo nije
rađeno u dostupnom vremenu; ostaje otvoreno kao smjer nastavka
(\S\ref{ZakljucakNastavak}).

\textbf{Naučena komponenta se ne isplati na razvojnom hardveru.}
Poglavlje~\ref{NaucenaKomponenta} mjeri, ne pretpostavlja, da mala naučena
mreža na GPU-u bez matričnih/tenzorskih jedinica nosi cijenu (2,7~ms po
okviru na 1080p) koja nije opravdana dobitkom kvalitete (0,10--0,20~dB).
Ovo nije neuspjeh implementacije nego izravno mjerenje razloga zašto
komercijalna rješenja (DLSS, FSR4/Redstone, PSSR) traže namjenski ML
hardver — pitanje koje je opis teme postavio, s brojkom kao odgovorom.

\section{Ograničenja mjerne postavke}

\textbf{Jedna primarna testna scena.} Sponza je geometrijski bogata (tanka
geometrija, stupovi, natpisi) ali svjetlosno jednostavna (jednostavan
forward shading, mipmapirane teksture s malo aliasinga). Posljedica je
mjerljiva u poglavljima~\ref{ImplUpscaler} i \ref{ImplFrameGen}: ni
upscaler ni generiranje okvira na ovoj sceni ne \emph{ubrzavaju} render
(cijena mehanizma je usporediva ili veća od cijene renderiranja lagane
scene), pa se dobitak izvještava kao kvaliteta, ne kao FPS pri istoj sceni.
Sekundarna proceduralna scena (pokretni objekti) koristi se posvuda gdje je
gibanje objekata, a ne samo kamere, relevantno (ablacije generiranja okvira
i naučene mješavine).

\textbf{Aproksimacija latencije.} Latencija se mjeri od uzorkovanja ulaza do
\emph{swapa} (predaje kompozitoru/monitoru), ne do stvarnog fotona na
zaslonu — razlika koju bi trebao mjeriti vanjski senzor (fotodioda ili
kamera visoke brzine), izvan opsega ovog rada. Usporedbe između načina
prikaza u ovom radu (s generiranjem, bez njega, s pacingom, bez njega) su
i dalje valjane jer sve mjere isti dodatak na istom stogu.

\textbf{Sub-pixel preciznost procjene gibanja.} Implementirani optical flow
(Poglavlje~\ref{ImplOpticalFlow}) radi cjelobrojnu pretragu na luminanciji
pola prikazne razlučivosti — najmanji razlučiv pomak je dva prikazna
piksela. Parabolička interpolacija oko minimuma kriterija podudaranja bi
smanjila kvantizacijsku pogrešku i ostaje kao poznat, nenapravljen korak.


\chapter{Zaključak i mogući nastavak}\label{Zakljucak}

\section{Doprinosi}

Ovaj rad je od nule, u C++20 i OpenGL~4.6 compute shaderima, implementirao
kompletan cjevovod temporalnog upscalinga i generiranja interpoliranih
okvira po uzoru na AMD FSR3: prostorni i temporalni upscaler s dilatacijom
vektora, disokluzijom, lockovima i adaptivnom rekonstrukcijom
(Poglavlje~\ref{ImplUpscaler}); piramidalnu procjenu gibanja iz slike
(Poglavlje~\ref{ImplOpticalFlow}); generiranje međuokvira s frame pacingom,
latencijom i kompozicijom korisničkog sučelja
(Poglavlje~\ref{ImplFrameGen}); i malu naučenu komponentu koja uči birati
između klasičnih kandidata (Poglavlje~\ref{NaucenaKomponenta}). Svaki dio je
mjeren naspram vlastite \emph{ground-truth} reference — pravog okvira i
pravog među-okvira koje samo vlastiti renderer može proizvesti — pa su
tvrdnje u ovom radu brojke, ne dojam.

Svaki korak cjevovoda mjerljivo nadmašuje prethodni (Tablica~\ref{table_ch10_sinteza}),
i cijeli mjerni paket (283 konfiguracije) reproduciran je i provjeren na dva
neovisna GPU-a dvaju različitih proizvođača (odjeljak~\ref{ProvjeraHardver}),
uz kvalitetu koja se poklapa preko te granice na $10^{-3}$--$10^{-1}$~dB.
Rad odgovara i na pitanje koje njegov opis izravno postavlja — isplati li se
naučena komponenta bez namjenskog ML hardvera — mjerenjem, ne pretpostavkom:
ne isplati se, i to je izmjeren, obrazložen odgovor
(Poglavlje~\ref{NaucenaKomponenta}), ne izostanak rezultata.

\section{Mogući nastavak}\label{ZakljucakNastavak}

\begin{itemize}
\item \textbf{Veći naučeni model uz namjenski ML hardver.} Analiza cijene
(Poglavlje~\ref{NaucenaKomponenta}) pokazuje da trend kvalitete s veličinom
mreže nije zasićen (najveći testirani model je i najbolji), a granica je
hardverska, ne algoritamska — isti model bi na GPU-u s matričnim/tenzorskim
jedinicama nosio bitno manji dio cijene okvira.
\item \textbf{Usporedba s Lucas--Kanade metodom procjene gibanja} (opcijski
dodatak iz izvornog plana, \S\ref{Rasprava}) — gradijentna metoda naspram
implementiranog block matchinga, po točnosti i cijeni.
\item \textbf{FLIP i LPIPS} kao dodatne perceptivne metrike, uz postojeći
PSNR/SSIM, posebno za slučajeve gdje se ta dva mjerenja razilaze
(Poglavlje~\ref{ImplUpscaler}, aliasirani sadržaj).
\item \textbf{Sub-pixel profinjenje procjene gibanja} (parabolička
interpolacija oko minimuma kriterija podudaranja), koje bi smanjilo
kvantizacijsku grešku opisanu u Poglavlju~\ref{Rasprava}.
\item \textbf{Provjera periodičnosti unutar block matchinga}, koja bi
riješila preostali izvor greške na periodičnim teksturama identificiran u
Poglavlju~\ref{ImplFrameGen} (težina toka uz game vektor).
\end{itemize}


%% Literatura — jezgra popisa već identificirana u PLAN.md; dopuniti prema
%% stvarno citiranim mjestima u tekstu.
\begin{thebibliography}{30}
\addcontentsline{toc}{chapter}{Literatura}

\bibitem{FleetWeiss}
{\sc D.~Fleet, Y.~Weiss}, {\em Optical Flow Estimation}, u: {\em Handbook of
Mathematical Models in Computer Vision}, Springer, 2006.

\bibitem{Karis}
{\sc B.~Karis}, {\em High Quality Temporal Supersampling}, ACM SIGGRAPH
Courses: Advances in Real-Time Rendering, 2014.

\bibitem{YangSurvey}
{\sc L.~Yang}, {\sc S.~Nehab}, {\sc P.~V.~Sander}, et al., {\em A Survey of
Temporal Antialiasing Techniques}, Computer Graphics Forum
{\bf 39}(2)(2020), 607--621.

\bibitem{RIFE}
{\sc Z.~Huang}, {\sc T.~Zhang}, {\sc W.~Heng}, {\sc B.~Shi}, {\sc S.~Zhou},
{\em Real-Time Intermediate Flow Estimation for Video Frame Interpolation},
ECCV 2022. (arXiv:2011.06294)

\bibitem{FILM}
{\sc F.~Reda}, {\sc J.~Kontkanen}, {\sc E.~Tabellion}, et al., {\em FILM:
Frame Interpolation for Large Motion}, ECCV 2022. (arXiv:2202.04901)

\bibitem{WangSSIM}
{\sc Z.~Wang}, {\sc A.~C.~Bovik}, {\sc H.~R.~Sheikh}, {\sc E.~P.~Simoncelli},
{\em Image Quality Assessment: From Error Visibility to Structural
Similarity}, IEEE Transactions on Image Processing {\bf 13}(4)(2004),
600--612.

\bibitem{FSR3GPUOpen}
AMD GPUOpen, {\em FidelityFX Super Resolution 3 — Frame interpolation,
Optical flow, Super resolution temporal}. Dostupno na
{\verb"https://gpuopen.com/manuals/fidelityfx_sdk/"}. (dokumentacija
korištena kao referenca, implementacija u ovom radu je vlastita — vidi
docs/reference/README.md u repozitoriju rada)

\bibitem{DLSSNvidia}
NVIDIA, {\em DLSS 2.0: A Big Leap In AI Rendering} i {\em DLSS 3: AI-Powered
Neural Graphics Innovations}. Dostupno na
{\verb"https://www.nvidia.com/en-us/geforce/news/"}.

\bibitem{XeSSIntel}
Intel, {\em XeSS 2 Whitepaper}. Dostupno na
{\verb"https://www.intel.com/content/www/us/en/developer/articles/technical/xess2-whitepaper.html"}.

\bibitem{ArmNSS}
Arm, {\em How Arm Neural Super Sampling Works}. Dostupno na
{\verb"https://developer.arm.com/community/arm-community-blogs/"}.

\bibitem{ReflexNvidia}
NVIDIA, {\em Reflex 2: Even Lower Latency Gameplay With Frame Warp}.
Dostupno na {\verb"https://www.nvidia.com/en-us/geforce/news/"}.

\end{thebibliography}


\chapter*{Sažetak}
\markboth{}{SAŽETAK}
\addcontentsline{toc}{chapter}{Sažetak}

Razlučivost i frekvencija osvježavanja zaslona rastu brže od
rasterizacijske snage GPU-a, pa igre sve češće rekonstruiraju sliku iz
manjeg broja stvarno izračunatih uzoraka — kroz prostor (upscaling) i kroz
vrijeme (generiranje interpoliranih okvira). Ovaj rad iz opisa tehnika AMD
FidelityFX Super Resolution 3 (FSR3), a ne iz izvornog koda, iznova
implementira cijeli klasičan cjevovod u C++20 i OpenGL~4.6 compute
shaderima: temporalni upscaler s dilatacijom vektora gibanja, odbacivanjem
povijesti po dubini, zaštitom tankih detalja (lockovi) i adaptivnom
rekonstrukcijom; piramidalnu procjenu gibanja iz slike (optical flow) za
sadržaj koji igrini vektori ne opisuju; i generiranje međuokvira s
frame pacingom, latencijom i kompozicijom korisničkog sučelja. Uz to je
implementirana i mala naučena komponenta koja uči birati između klasičnih
kandidata mješavine umjesto ručno napisane heuristike.

Ključna postavka je vlastiti renderer testne scene: kako se scena i kamera
u cijelosti kontroliraju, moguće je renderirati i \emph{ground truth} — puni
native okvir i pravi među-okvir na trenutku $t-0,5\Delta t$ — čime kvaliteta
postaje PSNR/SSIM brojka naspram poznate istine, a ne subjektivan dojam.
Svaki dio cjevovoda mjerljivo nadmašuje prethodni: puni upscaler dobiva do
$+4,56$~dB nad prostornim baseline-om, procjena gibanja uklanja do 87\%
gibanja koje igrini vektori ne vide, a generiranje okvira dobiva 7--10~dB
nad naivnim spajanjem dva okvira. Naučena komponenta mjerljivo podiže
kvalitetu (SSIM na svih osam mjerenih pogleda), ali se na GPU-u bez
matričnih/tenzorskih jedinica ne isplati — cijena joj nije opravdana
dobitkom, što je izmjeren, a ne pretpostavljen odgovor na to koliko daleko
ide analitički pristup bez namjenskog ML hardvera. Cijeli mjerni paket od
283 konfiguracije reproduciran je i provjeren na dva GPU-a različitih
proizvođača, s kvalitetom koja se poklapa preko te granice.

\section*{Ključne riječi}

upscaling, generiranje interpoliranih okvira, temporalni antialiasing,
procjena gibanja, optical flow, GPU compute shaderi, OpenGL, FSR3,
naučena mješavina


\chapter*{A UI System for Dynamic Resolution Upscaling and Artificial Frame Generation in Interactive Graphics}
\markboth{}{SUMMARY}

\section*{Summary}
\addcontentsline{toc}{chapter}{Summary}

Display resolution and refresh rate grow faster than GPU rasterisation
throughput, so games increasingly reconstruct the image from fewer actually
rendered samples, both spatially (upscaling) and temporally (interpolated
frame generation). This thesis reimplements, from the published description
of AMD FidelityFX Super Resolution 3 (FSR3) rather than from its source
code, the full classical pipeline in C++20 and OpenGL~4.6 compute shaders:
a temporal upscaler with motion-vector dilation, depth-based history
rejection, thin-detail protection (locks) and adaptive reconstruction; a
pyramidal image-space motion estimator (optical flow) for content that
engine motion vectors do not describe; and interpolated-frame generation
with frame pacing, latency, and UI compositing. A small learned component
is also implemented, which learns to pick among the classical blend
candidates instead of a hand-written heuristic.

The central design choice is a custom test-scene renderer: because the
scene and camera are fully controlled, it can render \emph{ground truth} —
a full native frame and a true interpolated frame at $t-0.5\Delta t$ — so
quality becomes a PSNR/SSIM number against a known answer, not a subjective
impression. Every pipeline stage measurably improves on the previous one:
the full upscaler gains up to $+4.56$~dB over the spatial baseline, motion
estimation recovers up to 87\% of motion that engine vectors miss, and
frame generation gains 7--10~dB over naively blending two frames. The
learned component measurably improves quality (higher SSIM on all eight
measured views) but does not pay for itself on a GPU without matrix/tensor
units — an answer obtained by measurement, not assumption, to how far an
analytical approach can go without dedicated ML hardware. The full
measurement suite (283 configurations) was reproduced and cross-checked on
two GPUs from different vendors, with quality metrics agreeing across that
boundary.

\section*{Keywords}

upscaling, interpolated frame generation, temporal antialiasing, motion
estimation, optical flow, GPU compute shaders, OpenGL, FSR3, learned blending


\chapter*{Životopis}
\addcontentsline{toc}{chapter}{Životopis}

% Životopis piše svaki kandidat OSOBNO, u prvom licu jednine: biografski
% podaci, sudjelovanje na natjecanjima, osvojene nagrade. Ovo su stvarni
% osobni podaci koje ja (Claude) ne mogu i ne smijem izmisliti — dolje su
% samo strukturni placeholderi za oba kandidata.

\section*{Ante Galić}
[PLACEHOLDER: životopis — Ante Galić, prvo lice jednine]

\section*{Ivo Janković}
[PLACEHOLDER: životopis — Ivo Janković, prvo lice jednine]

\end{document}
